\documentclass[11pt,a4paper]{article}
\usepackage[margin=2.0cm]{geometry}
\usepackage{mathptmx}
\usepackage{amsmath,amssymb}
\usepackage{xcolor}
\usepackage{enumitem}
\usepackage{fancyhdr,titlesec}
\usepackage[hidelinks]{hyperref}
\usepackage{listings}
\usepackage{fancyvrb}
\usepackage{booktabs}
\usepackage{array}
\usepackage{microtype}
\usepackage{fontspec}
\usepackage{newunicodechar}
\newunicodechar{❌}{{\fontspec{Adwaita Mono}❌}}
\setmonofont{Adwaita Mono}[Scale=MatchLowercase]
\emergencystretch=2em
\sloppy
\definecolor{ink}{HTML}{1F2A44}
\definecolor{accent}{HTML}{0F5EA8}
\definecolor{quote}{HTML}{5A6B7C}
\hypersetup{pdftitle={GA Deep Review - Round 1 (verbatim)},pdfauthor={Google Deep Review},colorlinks=true,urlcolor=accent,linkcolor=accent}
\titleformat{\section}{\Large\bfseries\color{ink}}{\thesection}{1em}{}
\titleformat{\subsection}{\large\bfseries\color{ink}}{\thesubsection}{0.8em}{}
\titleformat{\subsubsection}{\normalsize\bfseries\color{ink}}{\thesubsubsection}{0.8em}{}
\pagestyle{fancy}
\fancyhead[L]{\small\color{quote}GA Deep Review --- Round 1 (verbatim)}
\fancyhead[R]{\small\color{quote}\thepage}
\fancyfoot{}
\begin{document}
\begin{titlepage}
\vspace*{2.8cm}
{\color{ink}\Huge\bfseries Google Deep Review\\\bigskip}
{\color{ink}\Large Round 1 Report\\\bigskip}
{\color{quote}\large Verbatim reproduction}\vspace{1.5cm}

\begin{tabular}{@{}ll@{}}
\textbf{Document} & \texttt{REVISION\_REVIEW\_ROUND1.pdf} \\
\textbf{Source} & \texttt{GA\_DEEP\_REVIEW\_ROUND1.md} \\
\textbf{Origin} & Google Deep Review output, pasted by the author from their browser \\
\textbf{Date received} & November 2026 (first revision cycle) \\
\textbf{Paper reviewed} & \texttt{final\_submission/paper/paper.pdf} (first revision) \\
\bottomrule
\end{tabular}

\vspace{1.0cm}
\noindent\textbf{Faithfulness statement.} This document is a byte-order-preserving
reproduction of the reviewer's report. No reviewer wording has been added, removed,
paraphrased, corrected or re-ordered. The only changes are typographic: markdown
headings and emphasis are rendered as document structure, and markdown code spans
are rendered in a monospaced typeface. Where the review contains inline LaTeX-style
mathematics (e.g.\ \texttt{\$V(p) \textbackslash cap V\textsubscript{prot}(f)\$}),
it is typeset as inline mathematics exactly as written in the original. The original
markdown source file is included in the reproducibility package for byte-level diffing.

\vfill
\end{titlepage}
\section*{HIGH LEVEL SUMMARY}
\par\medskip\noindent\rule{\linewidth}{0.4pt}\medskip\par
\subsection*{Summary}
This paper presents a bounded state-graph architecture built on LangGraph for intent-based self-healing in Software-Defined Networks (SDN). To prevent probabilistic large language models (LLMs) from directly mutating network state or hallucinating recovery outcomes, the system confines the model to intent parsing and natural-language narration while interposing three deterministic Python gates: an evidence gate that evaluates latency sample distributions and telemetry freshness, a nine-rule safety validator that enforces topological and flow-protection constraints, and a post-action verification gate that re-measures live end-to-end latency to separate command installation from true outcome recovery. The artifact is evaluated across six scenarios and three ablation studies on a four-switch Mininet topology driven by a Ryu controller and a custom pure-Python switch, accompanied by a self-reported truth audit of the codebase and experimental claims.
\subsection*{Strengths}
\begin{itemize}[leftmargin=1.8em,itemsep=2pt]
\item The paper addresses an operationally vital challenge in network automation—separating command execution from verified SLA recovery and ensuring systems can safely decline to act when telemetry is stale, missing, or already compliant.
\item The deterministic three-gate workflow and six-state terminal taxonomy (\texttt{intent\_satisfied}, \texttt{partial\_improvement}, \texttt{recovery\_failed}, \texttt{guardrail\_blocked}, \texttt{abstained\_insufficient\_evidence}, and \texttt{lab\_offline}) provide a clean, auditable blueprint for bounding control loops and ensuring honest failure reporting.
\item The authors demonstrate commendable transparency by including a detailed limitations section and an accompanying truth audit (\texttt{AUDIT.md}) that candidly documents architectural boundaries, fixed code defects, and what the experimental testbed can and cannot substantiate.
\end{itemize}
\subsection*{Weaknesses}
\begin{itemize}[leftmargin=1.8em,itemsep=2pt]
\item The core architectural concept of confining LLMs to structured intent translation while delegating conflict detection, safety validation, and closed-loop assurance to deterministic verifiers is incremental relative to established IBN standards and recent verifier-gated LLM networking frameworks, which the paper neither cites nor compares against (Reference Paper: Intent-Based Networking - Concepts and Definitions, Alexander Clemm et al., 2022; Reference Paper: What do LLMs need to Synthesize Correct Router Configurations?, Rajdeep Mondal et al., 2023; Reference Paper: NetConfEval: Can LLMs Facilitate Network Configuration?, Changjie Wang et al., 2024; Reference Paper: LLNet: An Intent-Driven Approach to Instructing Softwarized Network Devices Using a Small Language Model, Antonino Angi et al., 2025; Reference Paper: NetIntent: Leveraging Large Language Models for End-to-End Intent-Based SDN Automation, Md. Kamrul Hossain and Walid Aljoby, 2025).
\item Similarly, the mathematical components for 75th-percentile sample-distribution SLA violation checking and composite effective-delay/utilization path ranking represent standard network monitoring and traffic engineering heuristics, and their empirical validation in Ablation 2 (Table VI) exhibits algebraic inconsistencies between the reported 10-sample means, phase peak latencies, and the 25\% violation rule of Equation (4) (Reference Paper: CONGA: Distributed Congestion-Aware Load Balancing for Datacenters, Mohammad Alizadeh et al., 2014; Reference Paper: Pingmesh: A Large-Scale System for Data Center Network Latency Measurement and Analysis, Chuanxiong Guo et al., 2015).
\item There is a major misalignment between the paper's framing around "Agentic AI" and "OpenFlow 1.3 SDN control" and the actual empirical evaluation: no LLM is ever invoked during the experiments (relying entirely on a deterministic keyword fallback across six static prompts rather than benchmarking natural-language parsing), the custom pure-Python switch never enforces the controller's priority-100 \texttt{flow-mod} rules in the data plane (instead forwarding every packet via the controller with \textasciitilde{}13.7 ms overhead), and link utilization remains at 0.000\% across all runs (testing static \texttt{netem} delay flags rather than real traffic saturation).
\item The formal specification contains internal inconsistencies with the evaluated system behavior—most notably, Safety Check 7 (\texttt{protected\_flows\_unmodified}) requires strict node-disjointness ($V(p) \cap V_{\text{prot}}(f) = \varnothing$) even though the video and protected voice flows inevitably share destination nodes $\{s_4, h_4\}$, Equation (7) defines post-action verification using a scalar inequality ($a \le \tau$) rather than the sample-distribution predicate of Equation (4) described in Section VIII-A, and the planner unconditionally drops the active path in Scenario S4 to install a candidate with higher configured delay ($12\text{ ms}$ vs. $8\text{ ms}$) instead of abstaining before execution.
\end{itemize}
\subsection*{Potential Issues And Suggestions}
\subsubsection*{[Introduction, Related Work \& Problem Scenario] (Pages: 1-2)}
\paragraph{Potential Mistakes and Improvements}
\begin{enumerate}[leftmargin=1.6em,itemsep=3pt]
\item \textbf{Alignment Between the "Agentic AI / LLM" Framing and the Non-LLM Evaluated System (Title, Abstract, Section I, Section II-D, Section X, and Section XI):}
\begin{itemize}[leftmargin=1.8em,itemsep=2pt]
\item \textbf{Potential Concern:} There is a notable disconnect between how the paper is framed in the Title (\emph{"Trustworthy Agentic AI for Intent-Based Self-Healing SDN"}), Abstract, Index Terms (\emph{"large language models, trustworthy AI"}), Section I, and Section XI—which center the motivation and conclusions on safely integrating Large Language Models (LLMs) into network control loops—and what is actually exercised in the empirical evaluation. Specifically, the Abstract states: \emph{"The central design decision is that the language model is not trusted with the decision to change the network. It is used only to parse intent, narrate evidence and write the final explanation,"} and Section XI (page 9, second paragraph) draws primary conclusions regarding LLM integration: \emph{"the value of an LLM in a control loop comes not from what it decides but from how much of the decision it is allowed to touch. Confining the model to interpretation and narration... was what allowed the system to be trusted."} However, Section II-D (page 2, left column) and Section X (page 9, fifth bullet) disclose that \emph{"In the configuration evaluated here there is no API key at all, so the agent runs its fully deterministic fallback and says so in the run record"} (which Section VI-E describes as a deterministic keyword parser). Consequently, all experimental scenarios (S1–S6) evaluate a deterministic keyword parser (\texttt{intent} node), a deterministic regex/rule guardrail, and deterministic Python functions inside a state machine over 6 fixed strings, leaving the typed-intent trust boundary between the probabilistic model $M$ and the deterministic components $D$ (Section IV-C)—specifically schema adherence, entity extraction errors on ambiguous natural-language phrasing, and LLM resistance to prompt injection beyond static keyword matching—untested. Furthermore, in the proposed architecture (Section IV-C and Section VI), even when an LLM is present, it is never involved in proposing candidate paths or network mutations; candidate path ranking and selection are performed entirely by a deterministic Python function (Eq. (6)), whereas the LLM only parses the initial string into a 5-tuple intent and formats narrative text.
\item \textbf{Impact \& Suggestion:} While the authors transparently note in Section II-D and Section X that omitting the LLM is a deliberate choice because the deterministic gate structure is the primary contribution, waiting until Section II-D to reveal that no language model is invoked in any of the experiments creates a mismatch with the front matter and conclusions. In addition, the statement in Section I that \emph{"the agent may propose; only verified evidence may decide"} followed by Steps 1–7 implies that the probabilistic model proposes candidate actions that the deterministic gates then filter, when in fact the deterministic planner proposes the actions. To prevent any impression of overclaiming, please consider either evaluating the pipeline with an active language model (such as a locally hosted open-weight model if external API keys are unavailable) across varied natural-language intent prompts and adversarial variations (e.g., \emph{NetConfEval} or \emph{IBNBench}) to empirically validate the $M \to D$ trust boundary, or explicitly stating upfront in the Abstract, Section I, and Section XI that the evaluated artifact exercises the deterministic gating architecture in isolation using a keyword-parser fallback without an active LLM, while clarifying in Section I that the deterministic planner—rather than the LLM—proposes candidate reroute actions.
\end{itemize}
\end{enumerate}
\begin{enumerate}[leftmargin=1.6em,itemsep=3pt]
\item \textbf{Omission of Closely Related Prior Art in Network Verification, Safe SDN Control, and LLM-Based Configuration (Section II-C, Section II-D, Section XI, and References):}
\begin{itemize}[leftmargin=1.8em,itemsep=2pt]
\item \textbf{Potential Concern:} Section II (\emph{Related Work}, page 2) and the References section (page 9) cite only seven references across four sub-areas (\texttt{[1]}–\texttt{[7]}), leaving two key subsections without supporting literature to contextualize the paper's claims:
\item \emph{Section II-C (Verification and safe control of network automation):} This subsection contains zero citations and frames the separation of \emph{command success} (installing a flow rule) from \emph{outcome success} (verifying data-plane performance recovery) without referencing the extensive foundational literature in SDN verification, safe updates, and active data-plane probing. Prior systems have long interposed deterministic verification gates between SDN controllers and the data plane to check invariants before rule installation (e.g., \emph{VeriFlow}, Khurshid et al., NSDI 2013; \emph{Header Space Analysis / NetPlumber}, Kazemian et al., NSDI 2012/2013), enforced safe consistent network updates (Reitblatt et al., SIGCOMM 2012), and performed active data-plane probing to verify real forwarding behavior and latency distributions (e.g., \emph{Automatic Test Packet Generation (ATPG)}, Zeng et al., CoNEXT 2012; \emph{Pingmesh}, Guo et al., SIGCOMM 2015).
\item \emph{Section II-D (LLMs for network operation):} This subsection cites only \emph{NetGPT} \texttt{[6]} (Meng et al., 2023), which focuses on packet-level traffic representation and generation rather than network configuration, intent translation, or SDN control. It omits directly comparable literature that combines natural-language intent translation or LLM agents with deterministic verifiers, compilers, and SDN controllers—such as \emph{LUMI} (Jacobs et al., USENIX ATC 2021), \emph{Verified Prompt Programming} (Mondal et al., HotNets 2023), and \emph{NetConfEval} (Wang et al., CoNEXT 2024).
\item \textbf{Impact \& Suggestion:} Expanding Section II-C, Section XI, and the References to cite established SDN invariant verification and active latency-probing frameworks would allow the authors to much more precisely differentiate their three-gate architecture (specifically, gating closed-loop intent recovery on pre-action capacity/protection checks and post-action end-to-end RTT sample distributions rather than static topological invariants). Similarly, supplementing \texttt{[6]} in Section II-D with literature specifically targeting verifier-gated LLM network configuration and intent translation would properly situate the work within current domain research.
\end{itemize}
\end{enumerate}
\begin{enumerate}[leftmargin=1.6em,itemsep=3pt]
\item \textbf{Inconsistent Factual Claim Regarding the Truth Audit (Section I, Pages 1–2 vs. Section XI, Page 9):}
\begin{itemize}[leftmargin=1.8em,itemsep=2pt]
\item \textbf{Definitive Flaw:} In the fourth contribution bullet at the end of Section I (bridging page 1, right column and page 2, left column), the text states that the accompanying truth audit \emph{"corrected \textbf{eighteen} factual claims and fixed sixteen code defects."} However, in Section XI (page 9, right column, third paragraph), the text states that the exact same truth audit \emph{"corrected \textbf{sixteen} factual claims and fixed sixteen code defects."}
\item \textbf{Suggestion:} Reconcile the number of corrected factual claims (\texttt{eighteen} vs. \texttt{sixteen}) between Section I (page 2) and Section XI (page 9) so that both sections report the exact count recorded in \texttt{final\_submission/AUDIT.md}.
\end{itemize}
\end{enumerate}
\begin{enumerate}[leftmargin=1.6em,itemsep=3pt]
\item \textbf{Logical Mismatch in Failure Mode Descriptions and Congestion Framing in Section III (Page 2):}
\begin{itemize}[leftmargin=1.8em,itemsep=2pt]
\item \textbf{Potential Concern:} Section III (\emph{Problem Scenario}, page 2, left and right columns) contains two internal consistency issues in how the motivating scenario and failure modes are described:
\end{itemize}
\end{enumerate}
1. \emph{Mismatch between the stated failure mode and its example:} In the second paragraph of the right column (page 2), the text introduces a failure mode as: \emph{"A second, subtler failure mode: the operator may act when they should not."} However, the sentence immediately explaining this claim describes the exact opposite behavior—\emph{failing to act when they should}: \emph{"If a single latency reading happens to fall just under the threshold, a system that trusts that reading concludes 'no problem' and does nothing—while the true behaviour is that the flow oscillates above and below the limit and the video is still dropping frames."} The next paragraph then states: \emph{"The system described in this paper is the demonstration of the first scenario against the second,"} which is ambiguous because Section III describes a single overarching problem scenario alongside distinct operator failure modes (acting without verifying whether a reroute worsened latency, acting unnecessarily on a healthy network as in S1/Ablation 3, and failing to act when a single sample or average masks an oscillating distribution breach as in Eq. (4)/Ablation 2).
2. \emph{Background-transfer congestion narrative vs. synthetic delay fault:} Section III motivates the problem with \emph{"an unrelated background transfer happens to traverse the same link the video uses. That link becomes congested, and the video's latency rises,"} and Figure 1 includes a \emph{"background source"} (\texttt{h2}). While the authors transparently state in Section I (Abstract), Section VII-B, and Section X that the traffic generator is not wired in (\texttt{0.000\%} utilization) and the fault is a synthetic \texttt{netem} delay (\texttt{+40 ms}), presenting the narrative in Section III and Figure 1 as a live background-transfer congestion scenario without noting there that \texttt{h2} is inactive in the evaluated scenario can momentarily mislead readers before they reach Section VII.
\begin{itemize}[leftmargin=1.8em,itemsep=2pt]
\item \textbf{Suggestion:} In Section III (page 2, right column), clearly separate the failure mode of \emph{acting when unwarranted} (Scenario S1 / Ablation 3) from \emph{failing to act when a single reading or mean masks an oscillating flow} (Ablation 2), and clarify the concluding sentence of Section III to state that the system addresses these distinct operational failure modes. Additionally, briefly note in Section III or the caption of Figure 1 that the congestion symptom on \texttt{s2-s4} is emulated via synthetic link-delay injection rather than live traffic from \texttt{h2}.
\end{itemize}
\paragraph{Minor Corrections and Typos}
\begin{itemize}[leftmargin=1.8em,itemsep=2pt]
\item \textbf{Page 1, Left Column, Abstract (lines 28–30):} The phrase \emph{"the datapath does not yet apply the controller's priority-100 rules, and the lab's reset does not restore the datapath's learned path"} uses implementation-specific details ("priority-100 rules", "learned path") without prior context in the Abstract. Consider rephrasing slightly to state that the custom Python switch forwards packets via controller learning state rather than enforcing installed flow-table priority rules.
\end{itemize}
\begin{itemize}[leftmargin=1.8em,itemsep=2pt]
\item \textbf{Page 2, Left Column, Section II-B:} Alongside the survey citation \texttt{[5]}, consider citing the standard IRTF/IETF definition of Intent-Based Networking (Clemm et al., RFC 9315, 2022) to ground the declarative intent terminology.
\end{itemize}
\subsubsection*{[System Model, Mathematical Formulation \& Guarded Workflow Algorithm] (Pages: 2-5)}
\paragraph{Potential Mistakes and Improvements}
\begin{enumerate}[leftmargin=1.6em,itemsep=3pt]
\item \textbf{Logical Contradiction in Safety Check 7 (\texttt{protected\_flows\_unmodified}) Regarding Shared Nodes and Endpoints (Section V-C, Page 4; Figure 1, Page 2):}
\begin{itemize}[leftmargin=1.8em,itemsep=2pt]
\item \textbf{Definitive Mathematical Inconsistency:} In Section V-C (Page 4), Safety Check 7 is formally defined as requiring both switch-edge disjointness and strict node disjointness between the proposed action's path $p$ and every protected flow $f$:
\end{itemize}
\end{enumerate}
$$E_p \cap E_{\text{prot}} = \varnothing \quad \text{and} \quad V(p) \cap V_{\text{prot}}(f) = \varnothing$$
where $V(p)$ is defined as the node set of the action's path and $V_{\text{prot}}(f)$ is the node set of a protected flow's current path. The accompanying prose reinforces this literally: \emph{"Check 7 is deliberately stricter than an edge test: a protected flow is only safe if the action leaves it on the same route, so sharing any node with a protected flow's current path counts as touching it."}
\begin{itemize}[leftmargin=1.8em,itemsep=2pt]
\item \textbf{Inconsistency with the Topology and Results:} According to Figure 1 (Page 2) and the experimental setup (Section VII-C, Page 6), the controlled video flow is $h_1 \to h_4$ and the protected voice flow is $h_3 \to h_4$, where both $h_3$ and $h_4$ attach solely to switch $s_4$. Consequently:
\end{itemize}
1. The protected voice flow $h_3 \to h_4$ traverses the path $(h_3, s_4, h_4)$, giving node set $V_{\text{prot}}(f) = \{h_3, s_4, h_4\}$ (or switch-node set $\{s_4\}$), while $E_{\text{prot}} = \varnothing$ since it traverses no inter-switch edges in $E_S$.
2. Every valid candidate path $p$ for the video flow $h_1 \to h_4$ must traverse switch $s_4$ to reach destination host $h_4$ (e.g., $(h_1, s_1, s_3, s_4, h_4)$), so $\{s_4, h_4\} \subseteq V(p)$ (or $\{s_4\} \subseteq V(p)$ if $V(p)$ tracks only switch nodes).
3. Therefore, $V(p) \cap V_{\text{prot}}(f) \supseteq \{s_4, h_4\} \neq \varnothing$ for \emph{every} possible path between $h_1$ and $h_4$.
\begin{itemize}[leftmargin=1.8em,itemsep=2pt]
\item \textbf{Potential Concern:} Under the mathematical definition of Check 7 as written, no candidate path for $h_1 \to h_4$ could ever pass safety validation when $h_3 \to h_4$ is a protected flow, yet Table III (Page 7) reports that all 9/9 safety checks pass in Scenarios S2, S3, and S4. This indicates a discrepancy between the formal mathematical specification in Section V-C and the actual predicate evaluated in code (for example, whether Check 7 in the implementation actually verifies that the action does not alter a protected flow's existing path in the snapshot, whether shared access switches and destination endpoints are exempted, whether $\mathcal{P} = \varnothing$ in the evaluated runs, or whether $V_{\text{prot}}(f)$ is only populated for flows with non-empty switch-edge sets $E_{\text{prot}} \neq \varnothing$). Clarifying the exact mathematical condition implemented in code is critical for reproducibility.
\end{itemize}
\begin{enumerate}[leftmargin=1.6em,itemsep=3pt]
\item \textbf{Inconsistency Between Lab-Reported Degraded Links ($C$) and Gate 1's "Suspect Links" in Candidate Ranking and Safety Validation (Sections IV-B, V-B, V-C, VI-B, Algorithm 1, and Section VIII-A, Pages 3–6):}
\begin{itemize}[leftmargin=1.8em,itemsep=2pt]
\item \textbf{Potential Concern (Reproducibility \& Formal Specification):} Across Sections IV-B, V-B, V-C, VI-B, Algorithm 1, and Section VIII-A, there are conflicting descriptions of how Gate 1 selects suspect links under zero link utilisation and how the excluded link set $C$ in Eq. (5) ($A(p) = [E_p \cap C = \varnothing]$), Eq. (6), and Safety Check 6 (\texttt{capacity\_respected}: $e \notin C$) is defined and populated:
\item In Section IV-B (Page 3), $C$ is part of the raw lab snapshot $\Sigma = (G, \mathcal{L}, \mathcal{F}, C)$ and is defined strictly as \emph{"the set $C$ of links the lab reports as degraded or not up."} Similarly, Algorithm 1 (Line 15, Page 5) states \texttt{diagnose: C <- degraded links} directly from the snapshot without referencing Gate 1's output.
\item In Section VI-B (Page 4), however, the prose states that Gate 1 (the evidence gate) \emph{"names a suspect link—either one whose utilisation is at or above the congestion threshold, or, failing that, the slowest link on the current flow's path... Naming a suspect is not decoration: the candidate ranking of Eq. (6) excludes paths that cross it, so an unnamed bottleneck would silently remove the congestion-awareness from the repair."}
\item Later, Section VIII-A (Page 6, Scenario S2) notes that \emph{"The evidence gate named s2-s4—the link the scenario was holding, and the slowest on the flow's current path—and then s1-s2 as suspect links; the ranking of Eq. (6) put every candidate that crossed the degraded link last"}, where \texttt{s1-s2} is an undegraded link on the initial path \texttt{s1-s2-s4} with \texttt{0.000\%} utilisation and normal configured delay.
\item \textbf{Impact on Execution Logic:} These descriptions create two clarity issues for independent reimplementation:
\end{itemize}
\end{enumerate}
1. Given that \texttt{s1-s2} has \texttt{0.000\%} utilisation and normal configured delay (\texttt{s2-s4} carries the injected \texttt{+40 ms} per-interface \texttt{netem} delay), it is unclear why \texttt{s1-s2} was also named as a suspect link alongside \texttt{s2-s4} when Section VI-B describes naming \emph{"a suspect link... the slowest link"} in the singular (e.g., whether Gate 1 ranks and returns all links on the current path \texttt{s1-s2-s4} in order of effective delay $d_e^{\text{eff}}$). Moreover, if Eq. (6) and Check 6 excluded all \emph{suspect} links named by Gate 1 rather than only the lab-reported degraded link (\texttt{s2-s4}), then naming \texttt{s1-s2} as a suspect link in Scenario S2 would also disqualify four-hop paths crossing \texttt{s1-s2} (such as \texttt{s1-s2-s3-s4}) under $\neg A(p)$ and Check 6 ($e \notin C$). Even more importantly, in Scenario S4 (where the network is healthy with $0.000\%$ utilisation and the lab-reported degraded set $C = \varnothing$, but the $20\text{ ms}$ target is unreachable), Gate 1's fallback rule names the slowest link on the current healthy path (\texttt{s1-s2-s4}) as a suspect link; if that healthy link were added to $C$, healthy links would fail Check 6 (\texttt{capacity\_respected}), whereas Algorithm 1 Line 16 (\texttt{drop the candidate equal to the path the flow is already on}) already forces the reroute off \texttt{s1-s2-s4} unconditionally.
2. Conversely, if Eq. (5) and Check 6 only use the lab-reported degraded link set $C$ from $\Sigma$, then Gate 1's fallback naming of the slowest link on the current path has no mathematical connection to Eq. (5) or Eq. (6). Explicitly distinguishing the lab-reported degraded link set $C_{\text{deg}}$ from Gate 1's diagnosed suspect link set $C_{\text{susp}}$ across Sections IV-B, V-B, V-C, VI-B, and Algorithm 1—and clarifying how Gate 1 selects multiple suspect links when utilisation is zero—would resolve this ambiguity.
\begin{enumerate}[leftmargin=1.6em,itemsep=3pt]
\item \textbf{Contradiction in the Boundary Condition $n < 2$ vs. "for two samples" in the Violation Predicate Explanation (Section V-A, Eq. 4, Page 3):}
\begin{itemize}[leftmargin=1.8em,itemsep=2pt]
\item \textbf{Definitive Flaw in Prose Explanation:} In Equation (4) (Page 3), the third and fourth rows evaluate the sample-distribution rule when $n \ge 2$ (which are mutually exclusive and exhaustive for all $n \ge 2$), while the fifth and sixth rows handle the fallback cases strictly when $n < 2$ (which for integer sample counts $n$ means $n \in \{0, 1\}$):
\end{itemize}
\end{enumerate}
$$\text{true, if } n < 2 \text{ and the peak} > \tau,$$
$$\text{true, if } n < 2, \text{ no peak, latency} > \tau,$$
However, the explanatory paragraph in the right column of Page 3 states:
> \emph{"Note what the last two lines do: when the lab reports fewer than two samples there is no distribution to speak of, so the rule falls back on the reported peak and then on the reported latency—which for two samples is the mean."}
\begin{itemize}[leftmargin=1.8em,itemsep=2pt]
\item \textbf{Potential Concern:} First, because the fallback branches only execute when $n < 2$ (\emph{"fewer than two samples"}), the fallback can never be reached when there are two samples ($n = 2$, which is evaluated by the distribution rule in rows 3–4 with $\gamma n = 0.5$). Either the threshold in the implementation is $n \le 2$ (requiring $n \ge 3$ samples to use the distribution rule, in which case Eq. 4 should be updated), or for $n = 1$ the reported latency is simply that single sample (or an aggregate scalar when individual sample arrays are absent). Second, the "last two lines" of Eq. (4) are actually the sixth line (\texttt{true, if n < 2, no peak, latency > \textbackslash{}tau}) and the seventh line (\texttt{false, otherwise}), whereas the two fallback rules for $n < 2$ are the fifth and sixth lines.
\end{itemize}
\begin{enumerate}[leftmargin=1.6em,itemsep=3pt]
\item \textbf{Incomplete Control-Flow Specification in Algorithm 1 and Figure 2 for \texttt{lab\_offline}, \texttt{guardrail\_blocked}, Risk-Tier Holds, and Early Exits (Sections V-C and VI, Pages 4–5):}
\begin{itemize}[leftmargin=1.8em,itemsep=2pt]
\item \textbf{Potential Concern (Algorithmic Completeness):} Section VI-A and the caption of Algorithm 1 (Page 5) state that \emph{"Algorithm 1 states the whole procedure, and the routing in it is the routing function of the implementation, not a paraphrase of it."} However, cross-checking Algorithm 1 and Figure 2 against Sections V-C, V-D, VI-A, and VI-E reveals several control-flow gaps:
\item \textbf{Unreachable lab (\texttt{lab\_offline} / \texttt{offline}):} Section VI-A (Page 4) explains that when the lab is unreachable, \emph{"The observation node fails, so the evidence gate has nothing to work with and the run takes the abstention branch—but the abstention is not allowed to overwrite the reason, and the recorded outcome stays \texttt{lab\_offline}."} In Algorithm 1 (Lines 5–13, Page 5), Line 5 notes \texttt{observe:} $\Sigma_0 \leftarrow \Sigma(I)$ \texttt{▷ fails if the lab is unreachable}, and Line 12 routes via \texttt{go to abstain} without showing how the terminal outcome $o \leftarrow \text{offline}$ (\texttt{lab\_offline}) is set or preserved over \texttt{abstain}. Furthermore, \texttt{lab\_offline} (as well as the specific Gate 3 verification outcomes \texttt{partial}, \texttt{failed}, and \texttt{unverified}) is not depicted among the terminal outcome nodes in Figure 2.
\item \textbf{Pre-graph guardrail (\texttt{guardrail\_blocked}):} Section VI-E (Page 5) describes a deterministic guardrail that screens the operator's request \emph{"Before the graph runs at all"} and terminates with \texttt{guardrail\_blocked} (evaluated in Scenario S5). Neither \texttt{guardrail\_blocked} nor this screening step appears in Algorithm 1's \texttt{Ensure:} clause, Algorithm 1's pseudocode, or Figure 2.
\item \textbf{Risk-tier hold vs. check failure (\texttt{held} vs. \texttt{rejected}):} Section V-C (Page 4) states that \emph{"A failing check also sets the risk tier, and a tier above the unattended threshold turns the verdict into a hold rather than an execution."} However, Section VI-A (Page 4) and Algorithm 1 (Lines 17–21) state that the safety gate routes to \texttt{hold} \emph{"for an approved action that does [require human approval]"}—i.e., only \emph{after} validator $V$ admits $p$ in Line 17 (which requires all nine checks to pass). If a failing check is what elevates the risk tier above the unattended threshold, an action with a failing check would not be admitted at Line 17 and would be rejected at Lines 18–19 (\texttt{if no p is admitted then o <- rejected}), never reaching the \texttt{else if V requires human approval then o <- held} branch at Line 20.
\item \textbf{Unassigned return variable $o$ on early exits and outcome label mapping:} In Algorithm 1, Lines 3 (\texttt{go to stop}), 8 (\texttt{go to satisfied}), and 12 (\texttt{go to abstain}) jump to undefined line labels without assigning $o$ or routing to \texttt{finish} (Line 27: \texttt{record o...; return o}), unlike Lines 19 and 21 (\texttt{o <- rejected; go to finish} and \texttt{o <- held; go to finish}). Additionally, \texttt{stop} appears in Line 3 and Figure 2, but is omitted from the \texttt{Ensure:} list of terminal outcomes in Algorithm 1, and the shorthand labels in Eq. (7)/Algorithm 1 (\texttt{satisfied}, \texttt{success}, \texttt{partial}, \texttt{failed}, \texttt{offline}) differ from the recorded outcome strings (\texttt{intent\_satisfied}, \texttt{partial\_improvement}, \texttt{recovery\_failed}, \texttt{lab\_offline}).
\end{itemize}
\end{enumerate}
\begin{enumerate}[leftmargin=1.6em,itemsep=3pt]
\item \textbf{Mathematical Notation and Scaling in the Network and Safety Formulations (Sections IV-A, V-C, and V-D, Pages 2–4):}
\begin{itemize}[leftmargin=1.8em,itemsep=2pt]
\item \textbf{Definitive Set-Notation Error for $E_p$ (Section IV-A, Page 2):} For a path $p = (v_0, \dots, v_k)$, its switch-edge set is written as $E_p = \{\{v_i, v_{i+1}\} \cap E_S\}$. Because $\{v_i, v_{i+1}\}$ is a single undirected edge (a 2-element set of vertices) whereas $E_S \subseteq E$ is a set of edges, intersecting a single edge with a set of edges is a type mismatch that yields $\varnothing$. This should be written as $E_p = \{\{v_i, v_{i+1}\} : 0 \le i < k\} \cap E_S$.
\item \textbf{Per-hop vs. path-wide datapath processing cost $c_{\text{dp}}$ (Section IV-A, Eq. 2, Pages 2–3):} Equation (2) models round-trip time as $\text{RTT}(p) \approx 2(D(p) + 2\delta_{\text{acc}}) + c_{\text{dp}}$, and Page 3 defines $c_{\text{dp}}$ as \emph{"the datapath's own per-hop processing cost,"} deriving $c_{\text{dp}} \approx 13.7\text{ ms}$ from the healthy 3-switch (2-switch-hop) baseline path \texttt{s1-s2-s4} where $D(p) = 8\text{ ms}$ and $\text{RTT} = 30.1\text{ ms}$ ($30.1 - 2(8 + 0.2) = 13.7\text{ ms}$). Because $c_{\text{dp}}$ enters Eq. (2) as an unscaled path-wide additive constant rather than being multiplied by the number of switch hops, referring to it as a \emph{"per-hop processing cost"} is mathematically inconsistent; it should either be described as the total round-trip datapath/controller processing overhead or scaled by hop count in Eq. (2).
\item \textbf{Formal Redundancy Between Safety Check 3 (\texttt{path\_well\_formed}) and Check 8 (\texttt{flow\_scope}) (Section V-C, Page 4):} Check 3 is defined as \emph{"$p_0 = s, p_k = t$, no repeated node, and every node of $p$ in $V$,"} while Check 8 is defined as \emph{"$p_0 = s$ and $p_k = t$, so the action may not insert or remove an endpoint."} As written mathematically, Check 8 is an exact subset of Check 3. If Check 8 validates a different field in the codebase (for example, verifying that the action's declared endpoints $(s_a, t_a)$ match the parsed intent's endpoints $(s_I, t_I)$ from Eq. 3, whereas Check 3 verifies that the path array $p$ starts and ends at the action's endpoints), the notation in Check 8 should distinguish intent endpoints from action endpoints.
\item \textbf{Protected-flow condition in Row 4 of the Verification Verdict and Unspecified Thresholds (Sections V-C, V-D, and VI-B, Page 4):} In Eq. (7), Row 3 returns \texttt{partial} if $a \le \tau \land \neg L \land Q$, and Row 4 returns \texttt{partial} if $b \text{ known and } a < b$. Because Row 4 does not include $\land Q$, any action that lowers latency ($a < b$) while harming a protected flow ($\neg Q$) falls through to Row 4 and evaluates to \texttt{partial} rather than \texttt{failed}. Clarifying whether Row 4 also requires $Q$ to hold—and stating the configured numerical values of the congestion threshold ($80\%$, shown only later in Fig. 6 on Page 9) and maximum switch-hop budget $h_{\max}$ (never stated) when defining Section V-C—would make the formal specification self-contained.
\end{itemize}
\end{enumerate}
\par\medskip\noindent\rule{\linewidth}{0.4pt}\medskip\par
\paragraph{Minor Corrections and Typos}
\begin{itemize}[leftmargin=1.8em,itemsep=2pt]
\item \textbf{Page 2, Figure 1:} Consider labeling the configured per-edge one-way delays ($d_e$) directly on the topology links in Figure 1 so readers can verify the path delay sums cited in Section VII-A ($8\text{ ms}$ for \texttt{s1-s2-s4}, $12\text{ ms}$ for \texttt{s1-s3-s4}, and $15\text{ ms}$ and $13\text{ ms}$ for the two four-hop paths).
\end{itemize}
\begin{itemize}[leftmargin=1.8em,itemsep=2pt]
\item \textbf{Page 4, Section V-C (Left Column, Paragraph 3), Section VI-C (Right Column, Paragraph 1), and Page 5, Algorithm 1 (Line 17):} The cross-reference \texttt{Eq. (V-C)} is used three times to refer to the nine safety checks in Section V-C, which are formatted as an enumerated list rather than a numbered equation. Either assign an equation number to the conjunction of the nine checks or change \texttt{Eq. (V-C)} to \texttt{Section V-C}.
\end{itemize}
\begin{itemize}[leftmargin=1.8em,itemsep=2pt]
\item \textbf{Page 4, Section V-D (Left Column, Paragraphs 1–2):}
\item In the paragraph preceding Eq. (7), calligraphic $\mathcal{Q}$ is used both for \emph{"the set of protected-flow evaluations"} and for \emph{"the conjunction over the protected-flow evaluations"}, whereas Eq. (7) uses italic $Q$ for the conjunction. Using $\mathcal{Q}$ for the set and $Q$ consistently for the conjunction in the prose avoids symbol overloading.
\item In the paragraph below Eq. (7), the prose states: \emph{"The second rule is a distinct outcome and not a special case of the third: the latency now meets the target but the flow has started losing packets, and the honest report is partial, not success. The third rule fires on improvement..."} Because Row 1 of Eq. (7) is \texttt{unverified} and Row 2 is \texttt{success}, the two \texttt{partial} rules being discussed are actually the \textbf{third} and \textbf{fourth} rules (rows) of Eq. (7).
\end{itemize}
\begin{itemize}[leftmargin=1.8em,itemsep=2pt]
\item \textbf{Page 4, Section VI-C (Right Column, Paragraph 2, Line 7):} Unrendered Markdown asterisks appear in the typeset text: \texttt{*evidence*} instead of italicized \emph{evidence}.
\end{itemize}
\subsubsection*{[Experimental Setup, Empirical Results \& Ablation Studies] (Pages: 5-9)}
\paragraph{Potential Mistakes and Improvements}
\begin{enumerate}[leftmargin=1.6em,itemsep=3pt]
\item \textbf{Methodological Validity and Mathematical Consistency of Ablation 2 (Section IX-C, Table VI, and Figure 4):}
\end{enumerate}
There are two substantive concerns regarding the setup and mathematical consistency of the violation-rule ablation presented in Section IX-C and Table VI:
\begin{itemize}[leftmargin=1.8em,itemsep=2pt]
\item \emph{Mathematical contradiction between Table VI outcomes (\texttt{no / yes}), Equation (4), Table II peaks, and Figure 4:} According to Section VII-C, each snapshot consists of $n = 10$ ICMP samples, and by Equation (4) (with $\gamma = 0.25$), a flow is treated as violated (\texttt{System = yes} / \texttt{true}) whenever $|\{i : s_i > \tau\}| \ge \gamma n = 0.25 \times 10 = 2.5$ (i.e., when at least $3$ out of $10$ samples exceed the target $\tau$). Conversely, a \texttt{System = no} (\texttt{false}) verdict requires that at most $2$ out of $10$ samples exceed $\tau$, and at least $8$ samples satisfy $s_i \le \tau$. In Table VI, 12 of the 14 rows report a sample mean strictly greater than $\tau$ while reporting \texttt{System / mean-only} as \texttt{no / yes}, creating direct algebraic contradictions with the maximum peak latencies reported in Table II and plotted in Figure 4:
\item For \texttt{baseline healthy} ($\tau = 30\text{ ms}$, $\text{Mean} = 32.0\text{ ms}$, \texttt{no / yes}), \texttt{healthy impossible target} ($\tau = 30\text{ ms}$, $\text{Mean} = 32.5\text{ ms}$, \texttt{no / yes}), and \texttt{guardrail blocked} ($\tau = 30\text{ ms}$, $\text{Mean} = 32.1\text{ ms}$, \texttt{no / yes}), the reported snapshot means in Table VI are identical to the maximum \texttt{Peak} values reported across all $50$ individual samples of those respective phases in Table II ($32.0\text{ ms}$ for S1, $32.5\text{ ms}$ for S4, and $32.1\text{ ms}$ for S5). It is mathematically impossible for a 10-sample arithmetic mean to equal the maximum individual sample unless all $10$ samples equal that peak—in which case $100\%$ ($10/10$) of the samples exceed $30\text{ ms}$ and the system rule must evaluate to \texttt{yes}, not \texttt{no}. Even if \texttt{Peak} in Table II denoted the maximum snapshot mean rather than the maximum individual sample, achieving a 10-sample mean of $32.0\text{--}32.5\text{ ms}$ when at least $8$ samples are $\le 30.0\text{ ms}$ requires the remaining $2$ samples to average at least $40.0\text{--}42.5\text{ ms}$, which contradicts Figure 4 where all healthy samples cluster tightly near $30\text{ ms}$.
\item Similarly, across the eight degraded rows (\texttt{congestion recovery} for S2 and \texttt{congestion partial} for S3), Table VI reports $\tau = 110\text{ ms}$, sample means ranging from $110.4\text{ ms}$ to $112.3\text{ ms}$, and verdicts of \texttt{no / yes}. For $10$ samples to achieve a mean of $112.3\text{ ms}$ while at least $8$ samples are $\le 110.0\text{ ms}$, the $2$ exceeding samples would have to average at least $121.5\text{ ms}$ (since $(8 \times 110.0 + 2 \times 121.5)/10 = 112.3$). Yet Table II states that across all $25$ snapshots ($250$ samples) of \texttt{S2 degraded} and all $5$ snapshots ($50$ samples) of \texttt{S3 degraded}, the maximum observed \texttt{Peak} was $112.2\text{ ms}$ (smaller than the $112.3\text{ ms}$ mean in Table VI) and $111.2\text{ ms}$, respectively (with standard deviations of only $0.9\text{ ms}$ and $0.7\text{ ms}$). Even at the absolute upper bound where $8$ samples equal $110.0\text{ ms}$ and $2$ samples equal the phase peak ($112.2\text{ ms}$ in S2 or $111.2\text{ ms}$ in S3), the maximum mathematically possible 10-sample mean that avoids triggering Equation (4) is $(8 \times 110.0 + 2 \times 112.2)/10 = 110.44\text{ ms}$ in S2 and $110.24\text{ ms}$ in S3. Thus, means of $110.7\text{--}112.3\text{ ms}$ in S2 and $110.4\text{--}111.1\text{ ms}$ in S3 are mathematically impossible alongside a \texttt{System = no} verdict. This is a definitive numerical inconsistency indicating that the script generating \texttt{ablations.csv} either inverted the boolean condition/inequality for the \texttt{System} column (e.g., requiring $>75\%$ rather than $\ge 25\%$ of samples to exceed $\tau$), or evaluated summary values rather than the 10-sample distributions.
\item \emph{Use of post-hoc synthetic latency targets disconnected from scenario intents:} According to Table I and Table II, the configured operator latency intents ($\tau$) are $50\text{ ms}$ for S1, S2, and S6, $30\text{ ms}$ for S3, $20\text{ ms}$ for S4, and no latency target at all (\texttt{–}) for S5 (\texttt{guardrail\_blocked}, where Section VIII-A further notes that the guardrail refused the request \emph{"before any telemetry was read"}). Under those actual scenario intents, Figure 4 shows that all individual ping samples lie more than $10\text{ ms}$ away from their respective operational targets, meaning the 25\%-sample rule and the mean-only rule produce identical verdicts during actual scenario execution. In Table VI, however, the \texttt{Target} column evaluates \texttt{baseline healthy} (S1), \texttt{healthy impossible target} (S4), and \texttt{guardrail blocked} (S5) against $\tau = 30\text{ ms}$, and \texttt{congestion recovery} (S2) and \texttt{congestion partial} (S3) against $\tau = 110\text{ ms}$. Because the empirical phase means are $\approx 29.7\text{ ms}$ (healthy) and $\approx 109.8\text{ ms}$ (degraded), setting the evaluation thresholds right at $30\text{ ms}$ and $110\text{ ms}$ places synthetic targets directly inside the empirical measurement noise. If synthetic thresholds ($\tau \in \{30, 110\}\text{ ms}$) were applied post-hoc to the recorded snapshots to illustrate boundary sensitivity, this methodology is a potential concern for clarity and should be explicitly explained and justified in Section IX-C and the caption of Table VI.
\end{itemize}
\begin{enumerate}[leftmargin=1.6em,itemsep=3pt]
\item \textbf{Cross-Table and Text-to-Table Discrepancies in Reported Latencies and Sample Notation (Abstract, Sections VIII–IX, Tables II–V, Table VII, and Figures 3–7):}
\end{enumerate}
A potential concern affecting clarity and reproducibility is that several key latency metrics are reported as different point values across the Abstract, prose, tables, and figures without explaining how individual trials or measurement sources were selected:
\begin{itemize}[leftmargin=1.8em,itemsep=2pt]
\item \emph{Scenario S2 repaired ("After") and "Before" latencies (\texttt{29.7 ms -> 109.8 ms, repaired to 38.7 ms} vs. \texttt{109.5} / \texttt{36.6 ms} vs. \texttt{41.6 ms}):} In the Abstract (page 1, left column) and Section VIII-A (page 6), the verified recovery in Scenario S2 is summarized as \texttt{(29.7 ms -> 109.8 ms, repaired to 38.7 ms)}, stating that \emph{"the verification gate measured 38.7 ms on average against the 50 ms intent."} Figure 5 lists the 5 individual S2 trials as having post-action latencies of \texttt{41.6 ms}, \texttt{40.5 ms}, \texttt{36.8 ms}, \texttt{38.1 ms}, and \texttt{36.6 ms}, whose arithmetic mean is indeed $38.72\text{ ms}$ ($\approx 38.7\text{ ms}$, matching the \texttt{36.6–41.6 ms} range in Table IV) and whose pre-action latencies (\texttt{112.0}, \texttt{112.9}, \texttt{111.6}, \texttt{112.2}, and \texttt{109.5} ms) average to $111.64\text{ ms}$ (whereas \texttt{29.7 ms} and \texttt{109.8 ms} in Table II are the pre-agent multi-snapshot phase means). However, Table III (\texttt{scenario\_summary.csv}) reports only the 5th trial (\texttt{Before = 109.5}, \texttt{After = 36.6}, the lowest-latency trial) as the single outcome for S2, whereas Table V (\texttt{ablations.csv}) and Figure 7 report only the 1st trial (\texttt{41.6 ms}, the highest-latency trial) as the system's congestion-aware latency alongside the 25-snapshot phase mean (\texttt{109.8 ms} from Table II) for the shortest-path ablation. All summary tables, figures, and the Abstract should consistently report the 5-trial mean (with range or standard deviation) or explicitly clarify which specific trial index or phase baseline mean is displayed.
\item \emph{Scenario S3 pre-action degraded latency (\texttt{110.4 ms} / Peak \texttt{111.2 ms} vs. \texttt{112.1 ms}):} Table II reports that across the 5 snapshots of S3, the mean latency is \texttt{110.4 ms} with a maximum \texttt{Peak} of \texttt{111.2 ms}. However, Section VIII-A, Table III, and Figure 5 report the agent's before-action latency in S3 as \texttt{112.1 ms}, which exceeds the phase peak in Table II. While Section VII-C notes that the agent's before/after measurement is taken independently of the 5 pre-run phase snapshots, explicitly noting this distinction in the captions of Table II and Table III—and relabeling column \texttt{n} in Table II (which denotes the number of 10-ping snapshots: \texttt{5} for S1, S3–S5 and \texttt{25} for S2) so it is not confused with the per-snapshot sample count $n = 10$ of Equation (4)—would resolve this apparent discrepancy.
\item \emph{Healthy baseline latency (\texttt{29.7 ms} vs. \texttt{29.5 ms} vs. \texttt{30.1 ms}):} The healthy video baseline latency on \texttt{s1-s2-s4} is reported as \texttt{29.7 ms} in Section VIII-A, Table II (S1 mean), and Table III (S4 Before), \texttt{29.5 ms, unchanged} in Table VII (\texttt{Policy: system}), and \texttt{30.1 ms} in the final row of Table VI (and earlier in Section IV-A). Since Table III lists \texttt{Before = –} for S1, the exact measurement source of the \texttt{29.5 ms} figure in Table VII (e.g., whether it is the agent's initial observation snapshot $\Sigma_0$ during S1) should be clarified.
\end{itemize}
\begin{enumerate}[leftmargin=1.6em,itemsep=3pt]
\item \textbf{Unexplained Second \texttt{always act} Row in Table VII and Undefined Phase \texttt{S1'} in Figures 3 and 4 (Sections IX-A and IX-D):}
\begin{itemize}[leftmargin=1.8em,itemsep=2pt]
\item \emph{Second \texttt{always act} row in Table VII:} Table VII (Ablation 3, page 8) contains two separate rows for the \texttt{always act} policy: the first shows \texttt{Outcome = recovery\_failed} with \texttt{Healthy flow = 29.7 -> 36.3 ms (+22\%)}, and the second shows \texttt{Outcome = --} with \texttt{Healthy flow = 36.3 -> 28.9 ms (-20\%)}. Section IX-D discusses only the first row (\texttt{29.7 ms to 36.3 ms}, corresponding to Scenario S4 where the healthy flow on \texttt{s1-s2-s4} is moved onto the slower \texttt{s1-s3-s4} path) and never mentions or explains the second row (\texttt{36.3 -> 28.9 ms (-20\%)}). It is a potential concern for reproducibility that the origin of this row is omitted—specifically whether it represents a subsequent always-act rollback/invocation starting from the already-rerouted \texttt{s1-s3-s4} path that oscillates the flow back to \texttt{s1-s2-s4} (and how that interacts with the third limitation in Section X stating that a single run does not reproduce the baseline after a reroute because switches retain learned rules), as well as why its outcome column is \texttt{--}.
\item \emph{Undefined category \texttt{S1'} in Figures 3 and 4:} The horizontal axes of both Figure 3 and Figure 4 plot six categories: \texttt{S1}, \texttt{S2 (5 trials)}, \texttt{S3}, \texttt{S4}, \texttt{S5}, and \texttt{S1'}, with the axis label reading \emph{"scenario (S1-S5 as in the scenario table)"}. However, \texttt{S1'} is never defined in Table I, Table II, the figure captions, or the body text. The authors should explicitly define \texttt{S1'} (e.g., whether it is a post-experiment healthy baseline check or corresponds to the second row of Table VII) and include its statistics in Table II if it is part of the experimental protocol.
\end{itemize}
\end{enumerate}
\begin{enumerate}[leftmargin=1.6em,itemsep=3pt]
\item \textbf{Exclusion of the Perturbed 6th Trial of Scenario S2 and Verification Gate Rule Attribution (Section VIII-A, Table IV, and Figure 5):}
\end{enumerate}
In Section VIII-A (page 6, right column), the authors state that S2 was repeated as 5 independent trials (all reporting \texttt{intent\_satisfied}), and then note: \emph{"in one further trial three of the ten post-action ping samples were inflated to about 90 ms by host scheduling, and because the gate counts every sample over the target as a breach—the same rule Eq. (4) uses—that trial was reported as a partial improvement. That is the gate behaving as specified rather than a failure, and it is the reason the trials are reported rather than a single run."} There are two potential concerns here:
\begin{itemize}[leftmargin=1.8em,itemsep=2pt]
\item \emph{Reporting completeness in Table IV and Figure 5:} If 6 trials were executed and the occurrence of host-scheduling perturbation is the stated motivation for reporting a multi-trial distribution, omitting the 6th trial from Table IV (\texttt{Trials: 5}, \texttt{Outcome(s): intent\_satisfied (5)}, \texttt{After (min–max): 36.6–41.6 ms}) and Figures 3–5 understates the empirical variance of the verification gate. Either all 6 trials (\texttt{intent\_satisfied (5), partial\_improvement (1)}) should be included in Table IV and Figures 3–5, or the protocol in Section VII-C should explicitly state the filtering criterion used to retain 5 trials.
\item \emph{Consistency between the prose explanation of Gate 3 and Equation (7):} The prose above states that the verification gate reported \texttt{partial\_improvement} because it \emph{"counts every sample over the target as a breach—the same rule Eq. (4) uses."} However, Equation (7) (Section V-D) defines the verification gate $\mathcal{V}$ using the scalar inequality $a \le \tau$ (where $a$ is the post-action flow latency), rather than invoking the 25\% sample-distribution predicate $\neg v(f^+, I)$ of Equation (4). While 3 samples at $\approx 90\text{ ms}$ and 7 samples at $\approx 37\text{ ms}$ would also push the 10-sample mean to $a \approx 52.9\text{ ms} > 50\text{ ms}$, the authors should clarify whether Gate 3 in the implementation evaluates $\neg v(f^+, I)$ (Equation 4) or the scalar mean $a \le \tau$ (Equation 7), especially given Section VII-D's statement that 27 unit tests check each equation of the paper directly against its implementation function.
\end{itemize}
\par\medskip\noindent\rule{\linewidth}{0.4pt}\medskip\par
\paragraph{Minor Corrections and Typos}
\begin{itemize}[leftmargin=1.8em,itemsep=2pt]
\item \textbf{Page 6, Section VII-C (Paragraph 1, Line 2):} \emph{"Each phase begins from a freshly started lab (Section VIII-A explains why that is necessary)..."} — Section VIII-A does not explain why the lab must be freshly started for each phase; that explanation appears in \textbf{Section X} (third bullet, Page 9: \emph{"The lab's reset does not restore the datapath's learned path... so the experiments restart the lab between scenarios instead"}). Update the cross-reference to \texttt{Section X}.
\end{itemize}
\begin{itemize}[leftmargin=1.8em,itemsep=2pt]
\item \textbf{Page 6, Section VIII-A (Paragraph 1, Line 1) \& Page 7, Table II:} \texttt{Tab. III} and \texttt{Tab. IV} are cited on Page 6 before \texttt{Table II} (Page 7) is ever cited in the text (\texttt{Table II} is never explicitly referenced in the body text of Sections VIII or IX). Ensure \texttt{Table II} is cited in the text and that table numbering follows the order of citation.
\end{itemize}
\subsubsection*{[Limitations, Truth Audit \& Conclusion] (Pages: 8-9)}
\paragraph{Potential Mistakes and Improvements}
\begin{enumerate}[leftmargin=1.6em,itemsep=3pt]
\item \textbf{Planner Logic Regression in Scenario S4 vs. the "Decides When Not to Act" Premise (Section VIII-A; Section X, Bullet 6; Section XI, Paragraph 1)}
\begin{itemize}[leftmargin=1.8em,itemsep=2pt]
\item \textbf{Potential Concern:} The paper's subtitle and central premise emphasize an agent \emph{"that Decides When Not to Act."} However, as disclosed in the sixth bullet of Section X (p. 9: \emph{"The planner compares configured and effective delays, not measured end-to-end latency. This is visible in scenario S4, where the agent selects a path that is slower in practice than the one it replaced"}), in Scenario S4 (healthy network with an unreachable 20 ms target, current measured RTT 29.7 ms on path \texttt{s1-s2-s4} with configured one-way switch delay $D(p) = 8\text{ ms}$), the planner unconditionally drops the current path (Algorithm 1, line 16) and selects the best remaining candidate (\texttt{s1-s3-s4} with $D(p) = 12\text{ ms}$), degrading the live flow's RTT from 29.7 ms to 36.3 ms. While Section X attributes this to comparing configured/effective delays rather than measured latency, according to Section VII-A even under the planner's own configured-delay metric ($D(p)$ in Eq. 1), the selected replacement path (\texttt{s1-s3-s4}, $12\text{ ms}$ one-way switch delay, implying $\text{RTT} \ge 24\text{ ms}$ before access and datapath overhead) is already strictly higher on paper than the active path (\texttt{s1-s2-s4}, $8\text{ ms}$) and cannot satisfy a $\tau = 20\text{ ms}$ RTT target. Having the planner execute a mutation that is guaranteed a priori to worsen delay—because Algorithm 1 (Line 16) unconditionally drops the current path from $\mathcal{P}_c$ without checking whether $\text{score}(p_{\text{candidate}}) < \text{score}(p_{\text{current}})$—and relying on Gate 3 (\texttt{verify}) to catch the regression after the live network has already been degraded contradicts the core goal of abstaining before mutating the network when action is unjustified.
\item \textbf{Suggested Improvement:} Update the candidate selection or safety validation logic (\texttt{diagnose} or \texttt{validate}) so that a candidate path $p_{\text{cand}}$ is only admitted if its score/projected delay improves upon the current active path $p_{\text{curr}}$ ($\text{score}(p_{\text{cand}}) < \text{score}(p_{\text{curr}})$) or has a plausible expectation of meeting the target $\tau$ (e.g., $2(D(p_{\text{cand}}) + 2\delta_{\text{acc}}) \le \tau$), and clarify in Sections VIII-A and X that even a planner using configured delays would have declined to act in S4 had it compared the candidate's configured score ($12\text{ ms}$) against the current path's configured score ($8\text{ ms}$).
\end{itemize}
\end{enumerate}
\begin{enumerate}[leftmargin=1.6em,itemsep=3pt]
\item \textbf{Unenforced OpenFlow Datapath Rules and Controller Packet-In Forwarding (Section VII-A; Section X, Bullets 2 \& 3; Section XI, Paragraph 1)}
\begin{itemize}[leftmargin=1.8em,itemsep=2pt]
\item \textbf{Potential Concern:} Section XI concludes that all six outcomes were produced on \emph{"a live OpenFlow lab"} matching the standard SDN architecture where a controller programs forwarding tables into switches. However, Section X (pp. 8–9, second and third bullets) discloses two critical infrastructure limitations in the custom pure-Python OpenFlow 1.3 switch: (a) \emph{"The datapath does not apply the controller's priority-100 rules. In the pure-Python datapath every data frame misses the flow table and is sent to the controller, which then forwards it; the effective path is therefore decided by the controller's learning state rather than by the table the controller installed,"} and (b) resetting the lab cannot clear learned switch rules without triggering broadcast storms across the ring topology, requiring full lab restarts between runs. Because SDN automation fundamentally relies on switches matching packets against installed flow-table entries in the data plane (and respecting rule priorities), forwarding every packet via controller \texttt{Packet-In} messages turns the testbed into controller-switched routing with a $\sim 13.7\text{ ms}$ per-packet processing overhead ($c_{\text{dp}}$, nearly double the $8\text{ ms}$ physical path delay) and host CPU scheduling ping spikes (up to $\approx 90\text{ ms}$) rather than standard OpenFlow data-plane forwarding.
\item \textbf{Suggested Improvement:} Fix the flow-table matching logic in the pure-Python switch so that priority-100 \texttt{OFPFlowMod} rules installed by the Ryu controller are matched directly in the switch datapath, or run the evaluation against standard Open vSwitch (OVS) in userspace mode (\texttt{datapath='user'} / \texttt{datapath\_type=netdev} in Mininet, which runs in unprivileged environments without kernel module dependencies) to verify actual data-plane rule enforcement and eliminate controller-round-trip scheduling jitter. Additionally, implementing Spanning Tree Protocol (STP) or explicit \texttt{(in\_port, eth\_dst)} / broadcast drop rules on the ring topology would eliminate broadcast storms during state resets and allow continuous multi-step runs.
\end{itemize}
\end{enumerate}
\begin{enumerate}[leftmargin=1.6em,itemsep=3pt]
\item \textbf{Untested Congestion-Utilization Logic Due to Zero Offered Load (Section VII-B; Section X, Bullet 1; Fig. 6)}
\begin{itemize}[leftmargin=1.8em,itemsep=2pt]
\item \textbf{Potential Concern:} Section VII-B, Section X (p. 8, first bullet), and Fig. 6 (p. 9) note that the injected fault is a static \texttt{netem} delay rather than offered traffic load because the existing traffic generator in the repository is not wired in, leaving measured link utilization ($u_e$) at \texttt{0.000\%} across all links and phases. Consequently, the utilization-based components of the mathematical model—specifically the tie-breaker $\frac{1}{100}\max_{e \in E_p} u_e$ in Eq. (6), the utilization threshold in Safety Check 6 (\texttt{capacity\_respected}), and the utilization branch of the Gate 1 suspect-link detector—are never exercised by live telemetry in the evaluation. Instead, Gate 1 and Gate 2 rely on the lab snapshot's scenario API explicitly flagging \texttt{s2-s4} in the degraded set $C$ ($e \notin C$) and setting $d_e^{\text{eff}}$. In a realistic SDN deployment, the physical network does not supply an external oracle flag $C$ stating which link has a delay fault; the system must infer degradation from live telemetry.
\item \textbf{Suggested Improvement:} Since Section VII-B notes that the traffic generator already exists in the repository, wire the background traffic generator (e.g., between \texttt{h2} and \texttt{h4}) into at least one scenario to test actual bandwidth saturation and exercise the congestion-detection and utilization-ranking branches on real switch port byte-counter telemetry, or clarify in Section X/XI how the system localizes a degraded link when the testbed does not explicitly pass the ground-truth degraded link set $C$ in the snapshot $\Sigma$.
\end{itemize}
\end{enumerate}
\begin{enumerate}[leftmargin=1.6em,itemsep=3pt]
\item \textbf{Missing Concrete Future Work Roadmap (Section XI)}
\begin{itemize}[leftmargin=1.8em,itemsep=2pt]
\item \textbf{Potential Concern:} While Section X candidly lists six major architectural and experimental limitations and Section XI provides a well-supported recommendation regarding shipping truth audits alongside research artifacts (given that several uncovered defects produced plausible numbers rather than obvious crashes), Section XI does not outline a concrete technical roadmap for resolving the identified system gaps.
\item \textbf{Suggested Improvement:} Add a brief paragraph in Section XI outlining concrete next steps for addressing the limitations identified in Section X (enforcing priority-100 rules in the switch datapath, integrating active background traffic loads, fixing the S4 planner comparison, exercising an active LLM across natural-language benchmarks, and evaluating multi-flow/scaled topologies such as the Internet Topology Zoo).
\end{itemize}
\end{enumerate}
\paragraph{Minor Corrections and Typos}
\begin{enumerate}[leftmargin=1.6em,itemsep=3pt]
\item \textbf{Bibliographic Errors in Reference [1] and Reference [7] (Page 9, References):}
\begin{itemize}[leftmargin=1.8em,itemsep=2pt]
\item \textbf{Reference [1] (OpenFlow):} The author list is given as \emph{"N. McKeown, T. Anderson, H. Aswani, V. Misra, A. Ranjan, D. Rajadeeswaran, J. Fernando, V. Guha, A. Wetherell, D. Wetherell, and R. Zhang"}. The actual authors of \emph{"OpenFlow: Enabling innovation in campus networks"} (\emph{ACM SIGCOMM CCR}, 2008) are \textbf{N. McKeown, T. Anderson, H. Balakrishnan, G. Parulkar, L. Peterson, J. Rexford, S. Shenker, and J. Turner}.
\item \textbf{Reference [7] (Mininet):} The author list is given as \emph{"B. Lantz, B. McGrew, and V. Paxson"}. The actual authors of \emph{"A network in a laptop: Rapid prototyping for software-defined networks"} (\emph{ACM HotNets}, 2010) are \textbf{B. Lantz, B. Heller, and N. McKeown}.
\end{itemize}
\end{enumerate}
\section*{CITATION \& REFERENCE VERIFICATION}
\par\medskip\noindent\rule{\linewidth}{0.4pt}\medskip\par
\subsection*{Citation Checker V2 Report}
\subsubsection*{Summary}
\begin{itemize}[leftmargin=1.8em,itemsep=2pt]
\item \textbf{Paper References Checked}: 7
\item \textbf{Verified}: 5
\item \textbf{Flagged (Likely Hallucinated)}: 2
\item \textbf{Unresolved}: 0
\end{itemize}
\subsubsection*{References Details}
\paragraph{❌ Reference 1: mckeown2008openflow}
\begin{itemize}[leftmargin=1.8em,itemsep=2pt]
\item \textbf{Final Verdict}: Likely Hallucinated
\item \textbf{Lookup Verdict}: Major Attribute Mismatch
\item \textbf{Search Verdict}: Hallucinated
\end{itemize}
\textbf{Extracted vs. Retrieved}:
\begin{itemize}[leftmargin=1.8em,itemsep=2pt]
\item \textbf{Authors}: Extracted: McKeown, N.; Anderson, T.; Aswani, H.; Misra, V.; Ranjan, A.; Rajadeeswaran, D.; Fernando, J.; Guha… — Retrieved: McKeown
\item \textbf{Journal}: Extracted: ACM SIGCOMM Computer Communication Review — Retrieved: —
\item \textbf{Pages}: Extracted: 69–76 — Retrieved: —
\item \textbf{Title}: Extracted: OpenFlow: Enabling innovation in campus networks — Retrieved: OpenFlow: Enabling Innovation in Campus Networks
\item \textbf{Volume}: Extracted: 38 — Retrieved: —
\item \textbf{Year}: Extracted: 2008 — Retrieved: —
\end{itemize}
\textbf{Search Rationale}:
> While the paper 'OpenFlow: Enabling innovation in campus networks' exists (published in ACM SIGCOMM Computer Communication Review, 2008), the author list is substantially hallucinated. The actual authors are Nick McKeown, Tom Anderson, Hari Balakrishnan, Guru Parulkar, Larry Peterson, Jennifer Rexford, Scott Shenker, and Jonathan Turner, whereas the citation invents nine fictitious co-authors (Aswani, Misra, Ranjan, Rajadeeswaran, Fernando, Guha, A. Wetherell, D. Wetherell, Zhang).
\textbf{Retrieved BibTeX}:
\begin{Verbatim}[fontsize=\small]
@misc{mckeown2008openflow,
  author = {McKeown},
  title = {{OpenFlow}: Enabling Innovation in Campus Networks},
}
\end{Verbatim}
***
\paragraph{❌ Reference 7: lantz2010network}
\begin{itemize}[leftmargin=1.8em,itemsep=2pt]
\item \textbf{Final Verdict}: Likely Hallucinated
\item \textbf{Lookup Verdict}: Major Attribute Mismatch
\item \textbf{Search Verdict}: Hallucinated
\end{itemize}
\textbf{Extracted vs. Retrieved}:
\begin{itemize}[leftmargin=1.8em,itemsep=2pt]
\item \textbf{Authors}: Extracted: Lantz, B.; McGrew, B.; Paxson, V. — Retrieved: Lantz
\item \textbf{Journal}: Extracted: Proceedings of the ACM SIGCOMM Workshop on Hot Topics in Networks (HotNets) — Retrieved: —
\item \textbf{Title}: Extracted: A network in a laptop: Rapid prototyping for software-defined networks — Retrieved: A Network in a Laptop: Rapid Prototyping for Software-Defined Networks
\item \textbf{Year}: Extracted: 2010 — Retrieved: —
\end{itemize}
\textbf{Search Rationale}:
> While the paper title 'A network in a laptop: Rapid prototyping for software-defined networks' and venue (ACM SIGCOMM HotNets 2010) are real, the authors listed in the citation ('McGrew, B. and Paxson, V.') are incorrect and hallucinated. The actual authors of the paper are Bob Lantz, Brandon Heller, and Nick McKeown.
\textbf{Retrieved BibTeX}:
\begin{Verbatim}[fontsize=\small]
@misc{lantz2010network,
  author = {Lantz},
  title = {A Network in a Laptop: Rapid Prototyping for {Software-Defined} Networks},
}
\end{Verbatim}
***
\end{document}
